\documentclass[12pt]{article}
\usepackage{amsmath, amssymb}
\usepackage{geometry}
\usepackage{array}
\usepackage{enumitem}
\usepackage{titlesec}
\usepackage[dvipsnames, table]{xcolor}
\usepackage{fancyhdr}
\usepackage{tcolorbox}
\usepackage{microtype}
\usepackage{tabularx}

\tcbuselibrary{skins, breakable}
\geometry{margin=1in, top=1in, bottom=1in}

\definecolor{primary} {RGB}{38,  50,  56}
\definecolor{accent}  {RGB}{0,  188, 212}
\definecolor{lightbg} {RGB}{224,247,250}
\definecolor{gold}    {RGB}{255,152,  0}
\definecolor{goldbg}  {RGB}{255,243,224}
\definecolor{softgray}{RGB}{245,246,247}
\definecolor{darktext}{RGB}{33,  33,  33}
\definecolor{green1}  {RGB}{0,  150, 136}
\definecolor{greenbg} {RGB}{224,242,241}

\pagestyle{fancy}
\fancyhf{}
\fancyhead[L]{\small\color{primary}\textbf{Operations Research}\;\textcolor{accent}{|}\;LP Formulation}
\fancyhead[R]{\small\color{darktext}Nadir Ali Khan \textbar\ MT-0123-052}
\renewcommand{\headrulewidth}{0pt}
\fancyhead[C]{\color{accent}\rule[-1ex]{\textwidth}{1.2pt}}
\fancyfoot[C]{\small\color{darktext}\thepage}

\titleformat{\section}{\large\bfseries\color{primary}}{}{0em}{}[\color{accent}\rule{\linewidth}{2pt}]
\titlespacing*{\section}{0pt}{20pt}{10pt}
\titleformat{\subsection}{\normalsize\bfseries\color{accent}}{}{0em}{}

\newtcolorbox{problembox}{enhanced, colback=lightbg, colframe=primary, arc=5pt, boxrule=2pt,
  left=12pt, right=12pt, top=8pt, bottom=8pt,
  title={\small\bfseries Problem Statement}, colbacktitle=primary, coltitle=white}

\newtcolorbox{notebox}{enhanced, borderline west={4pt}{0pt}{gold},
  colback=goldbg, colframe=white, arc=0pt, boxrule=0pt,
  left=12pt, right=8pt, top=6pt, bottom=6pt}

\newtcolorbox{defbox}{enhanced, borderline west={4pt}{0pt}{accent},
  colback=lightbg!60, colframe=white, arc=0pt, boxrule=0pt,
  left=12pt, right=8pt, top=6pt, bottom=6pt}

\newtcolorbox{stepbox}[1]{enhanced, colback=lightbg!60, colframe=accent, arc=4pt, boxrule=1.2pt,
  left=10pt, right=10pt, top=7pt, bottom=7pt,
  title={\small\bfseries #1}, colbacktitle=accent!80, coltitle=white,
  attach boxed title to top left={xshift=10pt, yshift=-\tcboxedtitleheight/2},
  boxed title style={arc=3pt, colframe=accent, colback=accent!70, boxrule=1pt}}

\newtcolorbox{answerbox}{enhanced, colback=greenbg, colframe=green1, arc=5pt, boxrule=1.5pt,
  left=12pt, right=12pt, top=10pt, bottom=10pt,
  title={\small\bfseries Optimal Solution}, colbacktitle=green1!85, coltitle=white,
  attach boxed title to top left={xshift=12pt, yshift=-\tcboxedtitleheight/2},
  boxed title style={arc=3pt, colframe=green1, colback=green1!70, boxrule=1pt}}

\newtcolorbox{iterbox}[1]{enhanced, colback=softgray, colframe=primary!40, arc=4pt, boxrule=1pt,
  left=10pt, right=10pt, top=6pt, bottom=6pt,
  title={\small\bfseries #1}, colbacktitle=primary!80, coltitle=white}

\newcolumntype{C}[1]{>{\centering\arraybackslash}p{#1}}

\begin{document}
\thispagestyle{empty}

\begin{tcolorbox}[enhanced, arc=0pt, boxrule=0pt, colback=primary, colframe=primary,
  left=16pt, right=16pt, top=18pt, bottom=14pt, borderline south={4pt}{0pt}{accent}]
  \centering
  {\huge\bfseries\color{white} Operations Research}\\[4pt]
  {\large\color{accent} Transportation Problem --- LP Formulation \& Solution}\\[8pt]
  \color{lightbg}\rule{0.5\textwidth}{0.4pt}\\[8pt]
  \begin{tabular}{rl}
    {\color{lightbg}\textbf{Name:}}       & {\color{white} Nadir Ali Khan} \\[3pt]
    {\color{lightbg}\textbf{Roll No:}}    & {\color{white} MT-0123-052} \\[3pt]
    {\color{lightbg}\textbf{Subject:}}    & {\color{white} Operations Research} \\[3pt]
    {\color{lightbg}\textbf{Question:}}   & {\color{white} Question No.\ 1} \\[3pt]
    {\color{lightbg}\textbf{Department:}} & {\color{white} BS Mathematics} \\
  \end{tabular}
\end{tcolorbox}

\vspace{14pt}

\begin{problembox}
\textbf{Question No.\ 1}\\[6pt]
\textbf{(a)} Define a transportation problem in operations research. Look at the network diagram given below and develop a linear programming problem.\\[6pt]
\textbf{(b)} Discuss balanced and unbalanced transportation problems with examples.
\end{problembox}

\vspace{10pt}

\begin{stepbox}{Transportation Cost Matrix (from Network Diagram)}
\vspace{4pt}
\begin{center}
\renewcommand{\arraystretch}{1.55}
\begin{tabular}{|c|C{2.2cm}|C{2.2cm}|C{2.2cm}|C{2.2cm}|}
  \hline
  \rowcolor{primary}
  {\color{white}\textbf{Source $\backslash$ Dest.}} &
  {\color{white}\textbf{$D_1$}} &
  {\color{white}\textbf{$D_2$}} &
  {\color{white}\textbf{$D_3$}} &
  {\color{white}\textbf{Supply $a_i$}} \\
  \hline
  \rowcolor{lightbg!50}
  \textbf{$S_1$} &
  $c_{11} = 3$ & $c_{12} = 4$ & $c_{13} = 2$ &
  \textbf{5} \\
  \hline
  \textbf{$S_2$} &
  $c_{21} = 2$ & $c_{22} = 2$ & $c_{23} = 3$ &
  \textbf{7} \\
  \hline
  \rowcolor{lightbg!50}
  \textbf{$S_3$} &
  $c_{31} = 3$ & $c_{32} = 1$ & $c_{33} = 3$ &
  \textbf{3} \\
  \hline
  \rowcolor{accent!20}
  \textbf{Demand $b_j$} & \textbf{7} & \textbf{3} & \textbf{5} &
  \textbf{15} \\
  \hline
\end{tabular}
\end{center}
\vspace{4pt}
{\small $c_{ij}$ = per-unit shipping cost from source $S_i$ to destination $D_j$.
Each cell entry is read from the arrow labels in the network diagram.}
\end{stepbox}

\vspace{10pt}

%=============================================================================
\section{Part (a): Definition and LP Formulation}

\subsection*{Definition of Transportation Problem}

\begin{defbox}
A \textbf{transportation problem} is a special class of linear programming problem concerned with distributing a homogeneous commodity from $m$ \emph{sources} (origins) to $n$ \emph{destinations} at \textbf{minimum total cost}.

Each source $S_i$ possesses a fixed supply $a_i$ and each destination $D_j$ requires a fixed demand $b_j$. The unit shipping cost from $S_i$ to $D_j$ is $c_{ij}$. The decision variables $x_{ij} \geq 0$ represent the quantity shipped on each route. A feasible solution must satisfy every supply and demand constraint exactly (for a balanced problem) or within bounds (for unbalanced cases).

Transportation problems arise in logistics, supply chain management, facility location, and production planning.
\end{defbox}

\vspace{8pt}

\subsection*{Decision Variables}

Let $x_{ij}$ = number of units transported from source $S_i$ to destination $D_j$\; ($i=1,2,3;\; j=1,2,3$).

\vspace{8pt}

\subsection*{LP Formulation}

\begin{stepbox}{Objective Function --- Minimise Total Transportation Cost}
\vspace{4pt}
\begin{align*}
\text{Minimise}\quad Z
  &= 3x_{11} + 4x_{12} + 2x_{13} \\
  &\quad + 2x_{21} + 2x_{22} + 3x_{23} \\
  &\quad + 3x_{31} + x_{32} + 3x_{33}
\end{align*}
\end{stepbox}

\vspace{8pt}

\begin{stepbox}{Supply and Demand Constraints}
\textbf{Supply constraints} (total shipped from each source $=$ its supply):
\begin{alignat}{2}
  x_{11} + x_{12} + x_{13} &= 5 \qquad &&(S_1) \notag\\
  x_{21} + x_{22} + x_{23} &= 7 \qquad &&(S_2) \notag\\
  x_{31} + x_{32} + x_{33} &= 3 \qquad &&(S_3) \notag
\end{alignat}

\textbf{Demand constraints} (total received at each destination $=$ its demand):
\begin{alignat}{2}
  x_{11} + x_{21} + x_{31} &= 7 \qquad &&(D_1) \notag\\
  x_{12} + x_{22} + x_{32} &= 3 \qquad &&(D_2) \notag\\
  x_{13} + x_{23} + x_{33} &= 5 \qquad &&(D_3) \notag
\end{alignat}

\textbf{Non-negativity:}\quad $x_{ij} \geq 0$ for all $i,j$.
\end{stepbox}

\vspace{8pt}

\begin{notebox}
\textbf{Balance Check:}\quad
Total Supply $= 5+7+3 = 15$,\quad Total Demand $= 7+3+5 = 15$.\\[3pt]
Supply $=$ Demand $\Rightarrow$ \textbf{Balanced Transportation Problem.} No dummy variable needed.
\end{notebox}

\vspace{10pt}

\subsection*{Solution: North-West Corner Method (Initial BFS)}

\begin{iterbox}{NWC Step-by-Step Allocations}
Start at cell $(1,1)$ and move right/down, always allocating $\min(\text{remaining supply, remaining demand})$.
\vspace{6pt}
\begin{center}
\renewcommand{\arraystretch}{1.35}
\begin{tabular}{|c|c|c|c|c|}
  \hline
  \rowcolor{primary}
  & {\color{white}$D_1$\;[7]} & {\color{white}$D_2$\;[3]} & {\color{white}$D_3$\;[5]} & {\color{white}\textbf{Supply}} \\
  \hline
  \rowcolor{lightbg!50}
  $S_1$\;[5] & \cellcolor{gold!35}$\mathbf{5}$\;(3) & 4 & 2 & 5 \\
  \hline
  $S_2$\;[7] & \cellcolor{gold!35}$\mathbf{2}$\;(2) & \cellcolor{gold!35}$\mathbf{3}$\;(2) & \cellcolor{gold!35}$\mathbf{2}$\;(3) & 7 \\
  \hline
  \rowcolor{lightbg!50}
  $S_3$\;[3] & 3 & 1 & \cellcolor{gold!35}$\mathbf{3}$\;(3) & 3 \\
  \hline
\end{tabular}
\end{center}
\vspace{4pt}
{\small Bold = allocation, parentheses = unit cost.}
\end{iterbox}

\vspace{6pt}

\begin{center}
\renewcommand{\arraystretch}{1.3}
\begin{tabular}{|C{2.6cm}|C{2.2cm}|C{1.8cm}|C{2.4cm}|}
  \hline
  \rowcolor{primary}
  {\color{white}\textbf{Route}} & {\color{white}\textbf{Units}} & {\color{white}\textbf{Cost}} & {\color{white}\textbf{Total}} \\
  \hline
  \rowcolor{lightbg!50} $S_1 \to D_1$ & 5 & 3 & 15 \\
  \hline
  $S_2 \to D_1$ & 2 & 2 & 4 \\
  \hline
  \rowcolor{lightbg!50} $S_2 \to D_2$ & 3 & 2 & 6 \\
  \hline
  $S_2 \to D_3$ & 2 & 3 & 6 \\
  \hline
  \rowcolor{lightbg!50} $S_3 \to D_3$ & 3 & 3 & 9 \\
  \hline
  \rowcolor{accent!25}
  \multicolumn{3}{|r|}{\textbf{Initial Cost}} & \textbf{40} \\
  \hline
\end{tabular}
\end{center}

\vspace{4pt}
Number of basic variables $= 5 = m+n-1 = 3+3-1$.\quad Non-degenerate. $\checkmark$

\vspace{10pt}

\subsection*{Optimality Test: MODI Method --- Iteration 1}

Set $u_1 = 0$. Compute duals $u_i,\,v_j$ via $c_{ij} = u_i + v_j$ for each basic cell:

\[
  v_1=3,\quad u_2=-1,\quad v_2=3,\quad v_3=4,\quad u_3=-1
\]

Reduced cost $\Delta_{ij} = c_{ij} - u_i - v_j$ for non-basic cells:

\begin{center}
\renewcommand{\arraystretch}{1.3}
\begin{tabular}{|c|c|c|c|c|}
  \hline
  \rowcolor{primary}
  {\color{white}\textbf{Cell}} & {\color{white}$c_{ij}$} &
  {\color{white}$u_i+v_j$} & {\color{white}$\Delta_{ij}$} & {\color{white}\textbf{Status}} \\
  \hline
  $(1,2)$ & 4 & $0+3=3$ & $+1$ & \\
  \hline
  \rowcolor{red!15}$(1,3)$ & 2 & $0+4=4$ & $\mathbf{-2}$ & \textbf{Enter (most negative)} \\
  \hline
  $(3,1)$ & 3 & $-1+3=2$ & $+1$ & \\
  \hline
  $(3,2)$ & 1 & $-1+3=2$ & $-1$ & \\
  \hline
\end{tabular}
\end{center}

\vspace{4pt}
Enter $(1,3)$.\quad
Loop:\;$(1,3)^+ \to (2,3)^- \to (2,1)^+ \to (1,1)^-$

$\theta = \min\{x_{23}=2,\;x_{11}=5\} = \mathbf{2}$

Updated: $x_{13}=2,\;x_{23}=0\text{ (leaves)},\;x_{21}=4,\;x_{11}=3$.\quad
Cost $= 40 - 2\times 2 = \mathbf{36}$

\vspace{10pt}

\subsection*{MODI Method --- Iteration 2}

New BFS: $x_{11}=3,\;x_{21}=4,\;x_{22}=3,\;x_{13}=2,\;x_{33}=3$

Duals ($u_1=0$): $v_1=3,\;v_3=2,\;u_2=-1,\;v_2=3,\;u_3=1$

\begin{center}
\renewcommand{\arraystretch}{1.3}
\begin{tabular}{|c|c|c|c|c|}
  \hline
  \rowcolor{primary}
  {\color{white}\textbf{Cell}} & {\color{white}$c_{ij}$} &
  {\color{white}$u_i+v_j$} & {\color{white}$\Delta_{ij}$} & {\color{white}\textbf{Status}} \\
  \hline
  $(1,2)$ & 4 & $0+3=3$ & $+1$ & \\
  \hline
  $(2,3)$ & 3 & $-1+2=1$ & $+2$ & \\
  \hline
  $(3,1)$ & 3 & $1+3=4$ & $-1$ & \\
  \hline
  \rowcolor{red!15}$(3,2)$ & 1 & $1+3=4$ & $\mathbf{-3}$ & \textbf{Enter (most negative)} \\
  \hline
\end{tabular}
\end{center}

Enter $(3,2)$.\quad
Loop:\;$(3,2)^+ \to (2,2)^- \to (2,1)^+ \to (1,1)^- \to (1,3)^+ \to (3,3)^-$

$\theta = \min\{x_{22}=3,\;x_{11}=3,\;x_{33}=3\} = \mathbf{3}$

Three cells tie; remove $(2,2)$. Cells $(1,1)$ and $(3,3)$ become $0$ (degenerate basic variables).

Updated allocations: $x_{32}=3,\;x_{21}=7,\;x_{13}=5$;\; $x_{11}=0,\;x_{33}=0$ (degenerate in basis).

Cost $= 36 - 3\times 3 = \mathbf{27}$

\vspace{10pt}

\subsection*{MODI Method --- Iteration 3 (Degenerate Pivot)}

BFS: $\{x_{11}=0,\;x_{21}=7,\;x_{13}=5,\;x_{32}=3,\;x_{33}=0\}$

Duals ($u_1=0$): $v_1=3,\;v_3=2,\;u_2=-1,\;u_3=1,\;v_2=0$

\begin{center}
\renewcommand{\arraystretch}{1.3}
\begin{tabular}{|c|c|c|c|c|}
  \hline
  \rowcolor{primary}
  {\color{white}\textbf{Cell}} & {\color{white}$c_{ij}$} &
  {\color{white}$u_i+v_j$} & {\color{white}$\Delta_{ij}$} & {\color{white}\textbf{Status}} \\
  \hline
  $(1,2)$ & 4 & $0+0=0$ & $+4$ & \\
  \hline
  $(2,2)$ & 2 & $-1+0=-1$ & $+3$ & \\
  \hline
  $(2,3)$ & 3 & $-1+2=1$ & $+2$ & \\
  \hline
  \rowcolor{red!15}$(3,1)$ & 3 & $1+3=4$ & $-1$ & Enter \\
  \hline
\end{tabular}
\end{center}

Enter $(3,1)$.\quad Loop:\;$(3,1)^+ \to (1,1)^- \to (1,3)^+ \to (3,3)^-$

$\theta = \min\{x_{11}=0,\;x_{33}=0\} = 0$\quad (pure degenerate pivot; no cost change)

Remove $(1,1)$; replace with $(3,1)=\varepsilon$. Cost remains $\mathbf{27}$.

\medskip
New BFS: $\{x_{31}=\varepsilon,\;x_{21}=7,\;x_{13}=5,\;x_{32}=3,\;x_{33}=0\}$

Recomputed duals: $u_1=0,\;u_2=0,\;u_3=1;\;v_1=2,\;v_2=0,\;v_3=2$

\begin{center}
\renewcommand{\arraystretch}{1.3}
\begin{tabular}{|c|c|c|c|c|}
  \hline
  \rowcolor{primary}
  {\color{white}\textbf{Cell}} & {\color{white}$c_{ij}$} &
  {\color{white}$u_i+v_j$} & {\color{white}$\Delta_{ij}$} & {\color{white}\textbf{Status}} \\
  \hline
  $(1,1)$ & 3 & $0+2=2$ & $+1$ & \\
  \hline
  $(1,2)$ & 4 & $0+0=0$ & $+4$ & \\
  \hline
  $(2,2)$ & 2 & $0+0=0$ & $+2$ & \\
  \hline
  $(2,3)$ & 3 & $0+2=2$ & $+1$ & \\
  \hline
\end{tabular}
\end{center}

All $\Delta_{ij} \geq 0$ $\Rightarrow$ \textbf{Optimal solution reached.}

\begin{answerbox}
\textbf{Optimal Allocation:}
\begin{center}
\renewcommand{\arraystretch}{1.35}
\begin{tabular}{|C{2.8cm}|C{2.4cm}|C{2cm}|C{2.6cm}|}
  \hline
  \rowcolor{green1!30}
  \textbf{Route} & \textbf{Units} & \textbf{Unit Cost} & \textbf{Total} \\
  \hline
  $S_1 \to D_3$ & 5 & 2 & 10 \\
  \hline
  $S_2 \to D_1$ & 7 & 2 & 14 \\
  \hline
  $S_3 \to D_2$ & 3 & 1 & 3 \\
  \hline
  \rowcolor{green1!40}
  \multicolumn{3}{|r|}{\Large\textbf{Minimum Transportation Cost}} & \Large\textbf{27} \\
  \hline
\end{tabular}
\end{center}
\vspace{4pt}
Verification:\quad $S_1$: $5=5$\checkmark\; $S_2$: $7=7$\checkmark\; $S_3$: $3=3$\checkmark\quad
$D_1$: $7=7$\checkmark\; $D_2$: $3=3$\checkmark\; $D_3$: $5=5$\checkmark
\end{answerbox}

%=============================================================================
\section{Part (b): Balanced and Unbalanced Transportation Problems}

\subsection*{1. Balanced Transportation Problem}

\begin{defbox}
A transportation problem is \textbf{balanced} when total supply equals total demand:
\[
  \sum_{i=1}^{m} a_i \;=\; \sum_{j=1}^{n} b_j
\]
Balanced problems are solved directly using standard methods (NWCM, VAM, LCM + MODI).
\end{defbox}

\textbf{Example} (the problem above):

\begin{center}
\renewcommand{\arraystretch}{1.3}
\begin{tabular}{|c|c|c|c|c|}
  \hline
  \rowcolor{primary}
  {\color{white}} & {\color{white}$D_1$} & {\color{white}$D_2$} & {\color{white}$D_3$} & {\color{white}Supply} \\
  \hline
  \rowcolor{lightbg!50} $S_1$ & 3 & 4 & 2 & 5 \\
  \hline
  $S_2$ & 2 & 2 & 3 & 7 \\
  \hline
  \rowcolor{lightbg!50} $S_3$ & 3 & 1 & 3 & 3 \\
  \hline
  \rowcolor{accent!20}\textbf{Demand} & 7 & 3 & 5 & \textbf{15 = 15} \\
  \hline
\end{tabular}
\end{center}

Total Supply $= 5+7+3 = 15\;=\;$ Total Demand $= 7+3+5 = 15$.\quad \textbf{Balanced.}

\vspace{10pt}

\subsection*{2. Unbalanced Transportation Problem}

\begin{defbox}
A transportation problem is \textbf{unbalanced} when total supply $\neq$ total demand. Two cases arise:

\begin{itemize}[leftmargin=*, itemsep=3pt]
  \item \textbf{Excess Supply} ($\sum a_i > \sum b_j$): Add a \textbf{dummy destination} $D_{n+1}$ with demand $= \sum a_i - \sum b_j$ and zero transportation costs $c_{i,n+1}=0$. Allocations to the dummy represent surplus (unsent) goods.
  \item \textbf{Excess Demand} ($\sum b_j > \sum a_i$): Add a \textbf{dummy source} $S_{m+1}$ with supply $= \sum b_j - \sum a_i$ and zero transportation costs $c_{m+1,j}=0$. Allocations from the dummy represent unmet (unsatisfied) demand.
\end{itemize}

After inserting the dummy, the problem becomes balanced and is solved by the standard algorithm.
\end{defbox}

\vspace{8pt}

\textbf{Example A --- Excess Supply ($\sum a_i > \sum b_j$):}

\begin{center}
\renewcommand{\arraystretch}{1.3}
\begin{tabular}{|c|c|c|c|c|}
  \hline
  \rowcolor{primary}
  {\color{white}} & {\color{white}$D_1$} & {\color{white}$D_2$} & {\color{white}$D_3$} & {\color{white}Supply} \\
  \hline
  \rowcolor{lightbg!50} $S_1$ & 5 & 3 & 6 & 20 \\
  \hline
  $S_2$ & 4 & 7 & 2 & 30 \\
  \hline
  \rowcolor{accent!20}\textbf{Demand} & 15 & 20 & 10 & \\
  \hline
\end{tabular}
\end{center}

Total Supply $= 50$,\quad Total Demand $= 45$,\quad Excess $= 5$.

Add dummy destination $D_4$ with demand $= 5$ and costs $= 0$:

\begin{center}
\renewcommand{\arraystretch}{1.3}
\begin{tabular}{|c|c|c|c|c|c|}
  \hline
  \rowcolor{primary}
  {\color{white}} & {\color{white}$D_1$} & {\color{white}$D_2$} & {\color{white}$D_3$} &
  {\color{white}$D_4$ (Dummy)} & {\color{white}Supply} \\
  \hline
  \rowcolor{lightbg!50} $S_1$ & 5 & 3 & 6 & 0 & 20 \\
  \hline
  $S_2$ & 4 & 7 & 2 & 0 & 30 \\
  \hline
  \rowcolor{accent!20}\textbf{Demand} & 15 & 20 & 10 & 5 & \textbf{50 = 50} \\
  \hline
\end{tabular}
\end{center}

Problem is now balanced and solvable. Units allocated to $D_4$ remain unsent at the origin.

\vspace{10pt}

\textbf{Example B --- Excess Demand ($\sum b_j > \sum a_i$):}

\begin{center}
\renewcommand{\arraystretch}{1.3}
\begin{tabular}{|c|c|c|c|c|}
  \hline
  \rowcolor{primary}
  {\color{white}} & {\color{white}$D_1$} & {\color{white}$D_2$} & {\color{white}$D_3$} & {\color{white}Supply} \\
  \hline
  \rowcolor{lightbg!50} $S_1$ & 2 & 8 & 5 & 10 \\
  \hline
  $S_2$ & 4 & 3 & 9 & 15 \\
  \hline
  \rowcolor{accent!20}\textbf{Demand} & 10 & 12 & 8 & \\
  \hline
\end{tabular}
\end{center}

Total Supply $= 25$,\quad Total Demand $= 30$,\quad Shortage $= 5$.

Add dummy source $S_3$ with supply $= 5$ and costs $= 0$:

\begin{center}
\renewcommand{\arraystretch}{1.3}
\begin{tabular}{|c|c|c|c|c|}
  \hline
  \rowcolor{primary}
  {\color{white}} & {\color{white}$D_1$} & {\color{white}$D_2$} & {\color{white}$D_3$} & {\color{white}Supply} \\
  \hline
  \rowcolor{lightbg!50} $S_1$ & 2 & 8 & 5 & 10 \\
  \hline
  $S_2$ & 4 & 3 & 9 & 15 \\
  \hline
  \rowcolor{lightbg!50} $S_3$ (Dummy) & 0 & 0 & 0 & 5 \\
  \hline
  \rowcolor{accent!20}\textbf{Demand} & 10 & 12 & 8 & \textbf{30 = 30} \\
  \hline
\end{tabular}
\end{center}

Problem is now balanced. Units allocated from $S_3$ represent demand that cannot be fulfilled.

\vspace{10pt}

\begin{notebox}
\textbf{Comparison Summary:}
\vspace{6pt}
\begin{center}
\renewcommand{\arraystretch}{1.4}
\begin{tabularx}{\linewidth}{|l|p{3.8cm}|X|}
  \hline
  \rowcolor{primary}
  {\color{white}\textbf{Type}} &
  {\color{white}\textbf{Condition}} &
  {\color{white}\textbf{Remedy}} \\
  \hline
  \rowcolor{lightbg!50}
  Balanced &
  $\displaystyle\sum a_i = \sum b_j$ &
  Solve directly using standard methods (NWCM, VAM, LCM + MODI). \\
  \hline
  Unbalanced\\(excess supply) &
  $\displaystyle\sum a_i > \sum b_j$ &
  Add a \textbf{dummy destination} $D_{n+1}$ with demand $= \sum a_i - \sum b_j$ and zero costs. \\
  \hline
  \rowcolor{lightbg!50}
  Unbalanced\\(excess demand) &
  $\displaystyle\sum b_j > \sum a_i$ &
  Add a \textbf{dummy source} $S_{m+1}$ with supply $= \sum b_j - \sum a_i$ and zero costs. \\
  \hline
\end{tabularx}
\end{center}
\end{notebox}

\vspace{20pt}
\noindent\color{accent}\rule{\linewidth}{1.5pt}
\begin{center}
  \small\color{darktext}\textit{End of Question No.\ 1\quad\textbullet\quad Operations Research\quad\textbullet\quad Nadir Ali Khan\quad\textbullet\quad MT-0123-052}
\end{center}

\end{document}
